两个等同两能级原子 \(|g\rangle,|r\rangle\) 被共振激光同时驱动，单原子Rabi频率为 \(\Omega>0\)。双激发态 \(|rr\rangle\) 因Rydberg相互作用具有额外能量 \(V>0\)。在旋转波近似的四态空间取
\[
H=\frac{\hbar\Omega}{2}(\sigma_x^{(1)}+\sigma_x^{(2)})+V|rr\rangle\langle rr|.
\]
起始 \(|gg\rangle\)，假设强阻塞 \(V\gg\hbar\Omega\)。
(1) 使用 \(|W\rangle=(|gr\rangle+|rg\rangle)/\sqrt2\) 与反对称态 \(|D\rangle=(|gr\rangle-|rg\rangle)/\sqrt2\)，写出完整的对称子空间哈密顿量与暗态解耦。
(2) 严格阻塞 \(V\to\infty\) 时解集体两态Rabi动力学，给最大单激发转移概率及产生纠缠W态的脉冲时间。
(3) 对大而有限的 \(V\) 用二阶消元求 \(|W\rangle\) 的光移及 \(|rr\rangle\) 的残余概率阶，解释误差为什么不随阻塞强度零次幂保持不变。
(4) 求理想 \(|W\rangle\) 的单原子约化态熵，检验其Bell纠缠与集体增强并非来自真正的双激发占据。
